\chapterbackmatter{Weinberg Exercises}

\begin{enumerate}
  \WeinbergExercise{1} Derive formulas \eqref{eq:31.2.3}--\eqref{eq:31.2.6} for
  components of the Einstein superfield.

  \WeinbergExercise{2} Suppose that supersymmetry is unbroken. Show how to calculate
  the amplitude for emission of a gravitino of very low energy in a
  general scattering process in terms of the amplitude for this process
  without the gravitino.

  \WeinbergExercise{3} Check that the supergravity action \eqref{eq:31.6.11} is
  invariant under the local supersymmetry transformations
  \eqref{eq:31.6.1}--\eqref{eq:31.6.6} to all orders in \(G\).

  \WeinbergExercise{4} Calculate the change of the generalized \(D\)-component
  \eqref{eq:31.6.40} under a general local supersymmetry transformation.

  \WeinbergExercise{5} Calculate the fermionic part of the Lagrangian density
  \eqref{eq:31.6.49}.

  \WeinbergExercise{6} Consider the theory of a single chiral scalar superfield
  \(\Phi\) interacting with supergravity, with modified Kahler
  potential \(d(\Phi,\Phi^*)=\Phi^*\Phi\) and superpotential
  \(f(\Phi)=M^2(\Phi+\beta)\), where \(M\) and \(\beta\) are constants.
  Find a value of \(\beta\) for which the classical field equations
  have a solution with flat spacetime. What is the value of \(\phi\)
  for this solution?
\end{enumerate}
